\documentclass[11pt,a4paper]{article}
\usepackage[utf8]{inputenc}
\usepackage[T1]{fontenc}
\usepackage{amsmath,amssymb}
\usepackage{booktabs}
\usepackage{graphicx}
\usepackage[margin=2.5cm]{geometry}
\usepackage{hyperref}
\usepackage{xcolor}

\hypersetup{
    colorlinks=true,
    linkcolor=blue!60!black,
    citecolor=blue!60!black,
    urlcolor=blue!60!black
}

\title{Multiple Scattering and Curvature Contributions\\
to the ALICE TRD Midpoint Residual}
\author{M.~Ivanov and the O2DistAI team\\[2pt]
{\normalsize JIRA: ATO-628, PHASE 0.4 --- TRD+TOF Alignment}}
\date{April 2026}

\begin{document}
\maketitle

%======================================================================
\section{Introduction}
%======================================================================

We estimate the multiple scattering (MS) and track curvature uncertainty
contributions to the midpoint residual in the ALICE Transition Radiation
Detector (TRD) for pions at $p = 1$, 2, 5, and 10~GeV/$c$.

The midpoint residual is defined as
\begin{equation}
\label{eq:midpoint}
\mathrm{Res}_i = d_i - \tfrac{1}{2}\bigl(d_{i-1} + d_{i+1}\bigr)\,,
\end{equation}
where $d_i$ denotes the \emph{residual of the cluster (tracklet) position
at TRD layer~$i$ with respect to the fitted curved track}.
The quantities $d_{i-1}$ and $d_{i+1}$ are the corresponding residuals
at the neighboring layers.

This estimator removes any linear component of the track
(straight-line propagation) but is sensitive to two effects:
\begin{enumerate}
\item \textbf{Multiple scattering} --- random angular deflections in the
      layer material produce a non-linear displacement that the midpoint
      formula cannot subtract.
\item \textbf{Track curvature mismatch} --- the residual retains the
      sagitta of the curved track; any uncertainty in the measured track
      momentum translates into a residual sagitta after the track
      fitter's curvature subtraction.
\end{enumerate}

%======================================================================
\section{Input parameters}
%======================================================================

\begin{table}[h]
\centering
\caption{Input parameters for the estimate.}
\label{tab:inputs}
\begin{tabular}{lll}
\toprule
Parameter & Value & Source \\
\midrule
$x/X_0$ per TRD layer & ${\approx}\,2.5\%$ (conservative) & Total TRD ${\sim}15\%\;/\;6$ layers \\
Inter-layer spacing $L$ & ${\approx}\,16$~cm & $\Delta r \approx 80$~cm, 6 layers \\
Layer thickness $t$ & ${\approx}\,10$~cm & Radiator + drift + readout \\
Magnetic field $B$ & 0.5~T & ALICE solenoid \\
TPC momentum resolution $\delta p/p$ & ${\sim}\,1$--$2\%$ & At $1$--$10$~GeV/$c$ \\
\bottomrule
\end{tabular}
\end{table}

\noindent
$L = 16$~cm is the adjacent-layer center-to-center spacing; the
outer-point chord (layer $i{-}1$ to layer $i{+}1$) is $2L = 32$~cm.
The formulas $s = L^2/(2R)$ and $s = (2L)^2/(8R)$ are equivalent.

The value $x/X_0 = 2.5\%$ is a conservative estimate based on the total
TRD radiation length of ${\sim}15\%$ divided equally among six layers.
Individual layer budgets vary from ${\sim}2.0\%$ to ${\sim}2.8\%$
depending on the chamber type and the distribution of material
(pad plane ${\sim}0.8\%$, radiator ${\sim}0.5\%$, gas ${\sim}0.2\%$,
support ${\sim}1\%$).

%======================================================================
\section{Multiple scattering contribution}
%======================================================================

\subsection{RMS scattering angle per layer}

The Highland formula gives the RMS projected scattering angle for a
particle traversing material of thickness $x/X_0$:
\begin{equation}
\label{eq:highland}
\theta_0 = \frac{13.6~\mathrm{MeV}}{\beta c\,p}\,
\sqrt{x/X_0}\;\bigl[1 + 0.038\,\ln(x/X_0)\bigr]\,.
\end{equation}
With $x/X_0 = 0.025$ and $\beta \approx 1$ for pions above 1~GeV/$c$:
\begin{equation}
\theta_0 \approx \frac{1.85~\mathrm{mrad}}{p~[\mathrm{GeV}/c]}\,.
\end{equation}

\subsection{MS contribution to the midpoint residual}

In the thin-scatterer model (measurement precedes scattering at each
layer), the midpoint residual~\eqref{eq:midpoint} picks up
contributions from:
\begin{enumerate}
\item \textbf{Scattering at layer~$i$} (dominant): the deflection
      $\theta_i$ shifts $d_{i+1}$ by $\theta_i L$ but does not affect
      $d_i$ or $d_{i-1}$.  Substituting into~\eqref{eq:midpoint}:
\begin{equation}
\mathrm{Res}_i^{(\mathrm{inter})} = -\tfrac{1}{2}\,\theta_i\,L
\qquad\Longrightarrow\qquad
\sigma_{\mathrm{inter}} = \tfrac{1}{2}\,\theta_0\,L\,.
\end{equation}

\item \textbf{Scattering within layer~$i$} (sub-dominant):
      scattering in the first half of the layer material
      displaces the particle at the measurement point:
\begin{equation}
\sigma_{\mathrm{intra}} =
\frac{\theta_{0,\mathrm{half}}\times t/2}{\sqrt{3}}
\approx \frac{36~\mu\mathrm{m}}{p~[\mathrm{GeV}/c]}\,.
\end{equation}

\item \textbf{Scattering at layers $i{\pm}1$}: their linear
      contributions cancel exactly in the midpoint formula (by
      construction).
\end{enumerate}

\noindent
Combined in quadrature:
\begin{equation}
\boxed{
\sigma_{\mathrm{MS}}(\mathrm{Res}_i)
= \sqrt{\sigma_{\mathrm{inter}}^2 + \sigma_{\mathrm{intra}}^2}
\approx \frac{152~\mu\mathrm{m}}{p~[\mathrm{GeV}/c]}
}
\end{equation}
The inter-layer term contributes 95\% of the variance.

%======================================================================
\section{Curvature (sagitta) contribution}
%======================================================================

\subsection{The sagitta}

A charged particle in a uniform magnetic field~$B$ follows a circular
arc of radius $R = p/(0.3\,B)$ (with $p$ in GeV/$c$, $B$ in~T, $R$
in~m).  The midpoint residual measures the sagitta~--- the deviation of
the middle point from the chord connecting the outer two:
\begin{equation}
s = \frac{L^2}{2R} = \frac{L^2 \times 0.3\,B}{2\,p}\,.
\end{equation}
With $L = 0.16$~m, $B = 0.5$~T:
\begin{equation}
s = \frac{1920~\mu\mathrm{m}}{p~[\mathrm{GeV}/c]}\,.
\end{equation}

\subsection{Sagitta uncertainty from momentum resolution}

The track fitter subtracts the best-estimate curvature from the cluster
positions.  What remains in the midpoint residual is the \emph{uncertainty}
in that subtraction:
\begin{equation}
\boxed{
\delta s = s \times \frac{\delta p}{p}
= \frac{1920}{p}\times\frac{\delta p}{p}~\mu\mathrm{m}
}
\end{equation}

\begin{table}[h]
\centering
\caption{Sagitta and its uncertainty from TPC momentum resolution.}
\label{tab:sagitta}
\begin{tabular}{ccccc}
\toprule
$p$ (GeV/$c$) & $s$ ($\mu$m) & $\delta s$ at 1\% & $\delta s$ at 1.5\% & $\delta s$ at 2\% \\
\midrule
1  & 1920 & 19.2 & 28.8 & 38.4 \\
2  &  960 &  9.6 & 14.4 & 19.2 \\
5  &  384 &  3.8 &  5.8 &  7.7 \\
10 &  192 &  1.9 &  2.9 &  3.8 \\
\bottomrule
\end{tabular}
\end{table}

\noindent
The full sagitta is large (mm scale) but the track fitter removes most
of it.  Only the residual from $\delta p/p$ uncertainty survives, which
at 1--2\% resolution is $\mathcal{O}(1$--$40)$~$\mu$m --- comparable to
or smaller than the MS contribution.

%======================================================================
\section{Combined results}
%======================================================================

All contributions are added in quadrature:
\begin{equation}
\sigma(\mathrm{Res}_i) = \sqrt{
\sigma_{\mathrm{MS}}^2 + \delta s^2 + \sigma_{\mathrm{intrinsic}}^2}\,,
\end{equation}
where $\sigma_{\mathrm{intrinsic}} = \sqrt{3/2}\;\sigma_{\mathrm{det}}
\approx 490~\mu$m for $\sigma_{\mathrm{det}} \approx 400~\mu$m
(the midpoint of three independent measurements with equal resolution
has variance $\tfrac{3}{2}\sigma_{\mathrm{det}}^2$).

\begin{table}[h]
\centering
\caption{Full breakdown of midpoint residual contributions
($\delta p/p = 1.5\%$).}
\label{tab:full}
\begin{tabular}{cccccccc}
\toprule
$p$ & $\theta_0$ & $\sigma_{\mathrm{MS}}$ & $s$ & $\delta s$ &
$\sigma_{\mathrm{intr}}$ & $\sigma_{\mathrm{total}}$ &
$\frac{\sigma_{\mathrm{MS}}\oplus\delta s}{\sigma_{\mathrm{intr}}}$ \\
(GeV/$c$) & ($\mu$rad) & ($\mu$m) & ($\mu$m) & ($\mu$m) &
($\mu$m) & ($\mu$m) & \\
\midrule
1  & 1850 & 152 & 1920 & 29 & 490 & 514 & 32\% \\
2  &  925 &  76 &  960 & 14 & 490 & 496 & 16\% \\
5  &  370 &  30 &  384 &  6 & 490 & 491 &  6\% \\
10 &  185 &  15 &  192 &  3 & 490 & 490 &  3\% \\
\bottomrule
\end{tabular}
\end{table}

\begin{table}[h]
\centering
\caption{Non-intrinsic (physics) contributions only: MS $\oplus$
curvature.}
\label{tab:physics}
\begin{tabular}{ccccc}
\toprule
$p$ (GeV/$c$) & $\sigma_{\mathrm{MS}}$ ($\mu$m)
& $\delta s$ ($\mu$m) & Combined ($\mu$m) &
Fraction of $\sigma_{\mathrm{intr}}$ \\
\midrule
1  & 152 & 29 & 155 & 32\% \\
2  &  76 & 14 &  77 & 16\% \\
5  &  30 &  6 &  31 &  6\% \\
10 &  15 &  3 &  15 &  3\% \\
\bottomrule
\end{tabular}
\end{table}

%======================================================================
\section{Interpretation}
%======================================================================

Multiple scattering dominates over curvature uncertainty at all momenta
in the 1--10~GeV/$c$ range.  The MS contribution is $5$--$10\times$
larger than~$\delta s$ because MS scales as $1/p$ while $\delta s$
scales as $(\delta p/p)/p$ --- at the 1--2\% momentum resolution level,
the curvature correction is already very good.

\begin{itemize}
\item At $p = 1$~GeV/$c$: combined physics contribution 155~$\mu$m
      (32\% of intrinsic).  Both MS and curvature are significant for
      alignment.

\item At $p = 2$~GeV/$c$: combined 77~$\mu$m (16\%).  MS dominates;
      curvature is a small correction.

\item At $p = 5$~GeV/$c$: combined 31~$\mu$m (6\%).  Both effects are
      sub-dominant.  This is the sweet spot for alignment extraction.

\item At $p = 10$~GeV/$c$: combined 15~$\mu$m (3\%).  Negligible
      relative to detector resolution; statistics become the limiting
      factor.
\end{itemize}

\medskip
\noindent\textbf{Key insight.}
The sagitta itself is large (1920~$\mu$m at 1~GeV/$c$, 192~$\mu$m at
10~GeV/$c$) --- but the track fitter removes it using the measured
momentum.  Only the \emph{uncertainty} in that removal ($\delta s =
s \times \delta p/p$) contributes to the midpoint residual.

\medskip
\noindent\textbf{Operational consequence.}
\emph{If structured residuals of 100--300~$\mu$m are observed at
$p > 5$~GeV/$c$, they are alignment or calibration effects, not
kinematic (MS or curvature) contamination.}
At these momenta the combined kinematic floor is below 31~$\mu$m.

\medskip
\noindent\textbf{Optimal momentum window for alignment extraction:}
$p \approx 3$--$10$~GeV/$c$, where MS\,+\,curvature combined are
below ${\sim}50~\mu$m (${<}10\%$ of intrinsic) while track statistics
remain adequate.  Below 3~GeV/$c$, MS contamination is significant;
above 10~GeV/$c$, statistics become sparse.

%======================================================================
\section{Notes}
%======================================================================

\begin{itemize}
\item The layer spacing $L = 16$~cm is the center-to-center distance
      between adjacent chambers (TRD radial extent ${\sim}80$~cm over
      6~layers).  The actual measurement-point spacing depends on
      whether the tracklet position is reconstructed at the anode plane
      or at the chamber midpoint.

\item The $x/X_0 = 2.5\%$ per layer is a conservative average.
      The Highland formula with average $x/X_0$ is accurate to
      ${\sim}10\%$.

\item For the $z$-direction midpoint residual, the same MS formula
      applies (scattering is isotropic in the small-angle limit).
      The curvature term is absent in $z$ (no bending in the
      solenoid field along the beam axis).

\item All numerical values are design-level estimates intended for
      order-of-magnitude guidance in TRD alignment weight
      parameterization.
\end{itemize}

\end{document}
